\documentclass[11pt]{article}

\usepackage[margin=1in]{geometry}
\usepackage{titlesec}
\usepackage{enumitem}
\usepackage{xcolor}
\usepackage{tabularx}
\usepackage{booktabs}
\usepackage{longtable}
\usepackage{hyperref}
\usepackage{fancyhdr}
\usepackage{parskip}

\usepackage{tcolorbox}
\usepackage{array}
\usepackage{multirow}
\setmainfont{Latin Modern Roman}

\definecolor{accent}{HTML}{1B4F72}
\definecolor{lightgray}{HTML}{F4F6F7}
\definecolor{warnyellow}{HTML}{FEF3CD}
\definecolor{warnyellowborder}{HTML}{D4A017}
\definecolor{successgreen}{HTML}{D4EDDA}
\definecolor{successborder}{HTML}{28A745}
\definecolor{dangerred}{HTML}{F8D7DA}
\definecolor{dangerborder}{HTML}{DC3545}
\definecolor{infoblue}{HTML}{D1ECF1}
\definecolor{infoborder}{HTML}{0C5460}

\tcbuselibrary{skins, breakable}

\newtcolorbox{notebox}[1][]{
  colback=infoblue, colframe=infoborder,
  fonttitle=\bfseries\small, title=#1,
  left=6pt, right=6pt, top=4pt, bottom=4pt,
  boxrule=0.6pt, arc=3pt
}

\newtcolorbox{warnbox}[1][]{
  colback=warnyellow, colframe=warnyellowborder,
  fonttitle=\bfseries\small, title=#1,
  left=6pt, right=6pt, top=4pt, bottom=4pt,
  boxrule=0.6pt, arc=3pt
}

\newtcolorbox{keybox}[1][]{
  colback=successgreen, colframe=successborder,
  fonttitle=\bfseries\small, title=#1,
  left=6pt, right=6pt, top=4pt, bottom=4pt,
  boxrule=0.6pt, arc=3pt
}

\titleformat{\section}{\large\bfseries\color{accent}}{\thesection}{0.6em}{}
\titleformat{\subsection}{\normalsize\bfseries\color{accent}}{\thesubsection}{0.6em}{}
\titleformat{\subsubsection}{\normalsize\itshape\color{accent}}{\thesubsubsection}{0.6em}{}

\hypersetup{
  colorlinks=true,
  linkcolor=accent,
  urlcolor=accent,
  citecolor=accent
}

\pagestyle{fancy}
\fancyhf{}
\fancyhead[L]{\small Alert \& Notification Microservice --- Proposal}
\fancyhead[R]{\small EMIS Development Team}
\fancyfoot[C]{\thepage}
\renewcommand{\headrulewidth}{0.4pt}

\setlist[itemize]{leftmargin=1.4em, itemsep=2pt}
\setlist[enumerate]{leftmargin=1.6em, itemsep=2pt}

\title{%
  \color{accent}\textbf{Alert \& Notification Microservice}\\[6pt]
  \large\normalfont\color{black}
  Proposal for a Standalone Alert Service in the EMIS Ecosystem
}
\author{EMIS Development Team}
\date{19 June 2026 \\ Branch: \texttt{be/feat/alert-service} \\ Status: Proposed}

\begin{document}

\maketitle
\thispagestyle{empty}
\vspace{0.4cm}

%---------------------------------------------------------------------------
\section*{Executive Summary}
%---------------------------------------------------------------------------

The Alert page (the officer ``Inbox'') is currently fetched directly from the Laravel backend,
mixing alert data with the rest of the HRM domain and giving future services no clean way to
raise notifications to officers. This proposal recommends building a dedicated, standalone
\textbf{Alert and Notification Microservice} in Node.js + Express with its own MySQL database.

The service is designed from day one to evolve from a polling-based data source (Phase~1)
into a fully event-driven publish model (Phase~2) without any change to the frontend contract.
It incorporates a hardened, layered security architecture covering JWT validation, rate limiting,
security headers, input sanitization, and database access controls.

The upstream dependency (currently the Laravel backend directly; WSO2 APIM in approximately
one month) is abstracted behind a single environment variable so the transition requires no
code change.

\newpage
\tableofcontents
\newpage

%---------------------------------------------------------------------------
\section{Background}
%---------------------------------------------------------------------------

The EMIS system manages the appointment and verification workflow for teachers and principals
across schools in Sri Lanka. When a profile is submitted it passes through a structured
approval chain involving verification officers, zonal education officers, and employer
representatives.

The Alert page is currently the only place in the system where officers can see what needs
their attention. It surfaces four categories of work:
\begin{itemize}
  \item Profiles waiting for first-level verification
  \item Profiles that were rejected and subsequently revised by the applicant
  \item Profiles waiting for employer confirmation
  \item Profiles that have been fully rejected
\end{itemize}
This page is critical to daily operations. Officers log in and check it to know what actions
they owe. Without it, the approval workflow has no visibility layer.

%---------------------------------------------------------------------------
\section{The Problem}
%---------------------------------------------------------------------------

\subsection{Alert data is coupled to the main Laravel backend}

\begin{itemize}
  \item \textbf{Alert data lives inside Laravel's database}, mixed with teacher profiles,
        institutions, timetables, and everything else. There is no domain separation.
  \item \textbf{Future services cannot contribute to the Alert page.} A future payroll or
        transfer-request service has no clean way to surface notifications to officers without
        writing directly into Laravel's database, defeating the purpose of having separate services.
  \item \textbf{The frontend code is inconsistent.} Alert data is currently fetched in three
        different ways across five different component files with no single source of truth.
  \item \textbf{No centralized security.} Each component calls the backend independently with
        no shared rate limiting, input validation, or audit logging for alert traffic.
\end{itemize}

\subsection{This does not scale}

As EMIS grows, each new service will generate its own events that officers need to see. Without
a dedicated alert service, every new service must couple to Laravel or the frontend must call
multiple backends to assemble a single notification view.

%---------------------------------------------------------------------------
\section{The Proposal}
%---------------------------------------------------------------------------

We propose building a \textbf{dedicated Alert and Notification Microservice}: a standalone
Node.js + Express backend with its own MySQL database and a hardened security stack.

\subsection{What it is}
\begin{itemize}
  \item A standalone HTTP service that owns all alert and notification data in its own database
  \item The service that powers the Alert page (Inbox) in the EMIS frontend
  \item The common ground for alerts across all current and future EMIS microservices
  \item A security-hardened service with layered defences (see Section~\ref{sec:security})
\end{itemize}

\subsection{What it is not}
\begin{itemize}
  \item Not a replacement for the main Laravel backend
  \item Not a messaging system --- no email or SMS in Phase~1
  \item Not a real-time push system --- no WebSockets in Phase~1
  \item Not a pass-through proxy --- it owns its own database from day one
\end{itemize}

%---------------------------------------------------------------------------
\section{Main Purpose}
%---------------------------------------------------------------------------

\begin{keybox}[Core purpose]
The alert service is the single place in the EMIS ecosystem where any service can publish
an alert, and any authorized user can read what needs their attention.
\end{keybox}

Today that means teacher and principal appointment workflow alerts for HR officers.
Tomorrow that means any event from any service --- a payroll discrepancy, a transfer
request, a document expiry --- flows through the same service and appears in the same
officer inbox.

The service gives EMIS a \textbf{common notification layer} independent of any single backend
technology. A Python service, a Go service, or a second Node service can all speak to one place.

%---------------------------------------------------------------------------
\section{Current Infrastructure Note}
%---------------------------------------------------------------------------

\begin{warnbox}[Important: APIM not yet in production]
WSO2 API Manager (APIM) is \textbf{not currently deployed} in the EMIS environment.
It is expected to be available in approximately one month. This affects the upstream
URL the alert service calls in Phase~1.
\end{warnbox}

The upstream dependency is abstracted behind a single environment variable
\texttt{UPSTREAM\_BASE\_URL}. Today it points directly to the Laravel backend;
when APIM is deployed, updating this one variable and redeploying is all that is
required. No code changes are needed.

\begin{table}[h]
\small
\begin{tabularx}{\textwidth}{l X}
\toprule
\textbf{Period} & \textbf{\texttt{UPSTREAM\_BASE\_URL} value} \\
\midrule
Now (no APIM) & \texttt{https://backend.app-uat.emis.moe.gov.lk/api} \\
In $\sim$1 month (with APIM) & \texttt{https://services.app-uat.emis.moe.gov.lk/hrm/1.0.0} \\
\bottomrule
\end{tabularx}
\end{table}

%---------------------------------------------------------------------------
\section{System Architecture}
%---------------------------------------------------------------------------

\subsection{Why a microservice}

The alert service is standalone --- not a module inside Laravel, not a proxy wrapper ---
because it \textbf{owns its own database}. That is what makes it independently deployable and
the common ground for all future EMIS services.

\begin{table}[h]
\small
\begin{tabularx}{\textwidth}{X X}
\toprule
\textbf{Problem with current setup} & \textbf{What the microservice solves} \\
\midrule
Alert data lives in Laravel's DB, coupled to the HRM domain
  & Alert service owns its own DB, fully decoupled \\
Future services cannot publish alerts without touching Laravel
  & Any service can publish to the alert service's API \\
Frontend fetches alerts from different sources inconsistently
  & One service, one contract, one URL \\
Cannot scale alert reads independently
  & Alert service scales on its own \\
No centralized security layer for alert traffic
  & Hardened middleware stack on every request \\
\bottomrule
\end{tabularx}
\end{table}

\subsection{System context --- current state (no APIM)}

\begin{quote}
\small\ttfamily
Browser (React SPA) \\
\hspace*{1em}Alert page $\rightarrow$ alertService.js $\rightarrow$ Alert Microservice URL \\[4pt]
$\downarrow$ HTTPS\quad Authorization: Bearer \textless token\textgreater \\[4pt]
Alert Microservice (Node.js + Express) \\
\hspace*{1em}Security stack: Helmet $\rightarrow$ CORS $\rightarrow$ Rate Limit $\rightarrow$ JWT Validate $\rightarrow$ Input Validate \\
\hspace*{1em}Owns: Alert MySQL DB (\texttt{alerts}, \texttt{alert\_counts}, \texttt{service\_keys}) \\
\hspace*{1em}Phase 1: sync middleware reads from Laravel directly \\[4pt]
$\downarrow$ HTTPS\quad Authorization: Bearer \textless token\textgreater \\[4pt]
Laravel Backend (FrankenPHP) \\
\hspace*{1em}JwtGuard validates token $\rightarrow$ returns alert data from its own DB
\end{quote}

\subsection{System context --- future state (with APIM in $\sim$1 month)}

\begin{quote}
\small\ttfamily
Browser (React SPA) \\[4pt]
$\downarrow$ \\[4pt]
Alert Microservice (Node.js + Express) \\[4pt]
$\downarrow$ HTTPS\quad Authorization: Bearer \textless token\textgreater \\[4pt]
WSO2 API Manager (APIM) \\
\hspace*{1em}Validates token $\rightarrow$ strips /hrm/1.0.0 $\rightarrow$ forwards to Laravel \\[4pt]
$\downarrow$ HTTP \\[4pt]
Laravel Backend (FrankenPHP) \\
\hspace*{1em}JwtGuard validates token $\rightarrow$ returns alert data from its own DB
\end{quote}

\begin{notebox}[Transition]
Switching from the current state to the APIM state requires changing \texttt{UPSTREAM\_BASE\_URL}
in the environment and redeploying. Zero code changes.
\end{notebox}

\subsection{Service responsibilities}

\textbf{In scope (Phase 1):}
\begin{itemize}
  \item Own and manage the alert MySQL database
  \item Expose a hardened HTTP API for reading alert data (counts + paginated lists)
  \item Validate every incoming JWT independently --- not relying solely on the upstream backend
  \item Rate-limit all endpoints to protect upstream services and the DB
  \item Validate and sanitize all input before it reaches the DB or upstream calls
  \item Phase 1: populate the DB by reading from the upstream backend via an internal sync middleware
  \item Normalize all data into a consistent camelCase shape before storing and serving
  \item Return consistent, information-safe error responses --- no stack traces, no internal details
\end{itemize}

\textbf{Out of scope (Phase 1):}
\begin{itemize}
  \item Accepting alert events published by producers (Phase~2)
  \item WebSocket / Server-Sent Events for real-time updates (future)
  \item Performing workflow actions (verify, confirm, reject) --- these remain direct
        frontend-to-backend calls
  \item Email or SMS notifications (future)
\end{itemize}

%---------------------------------------------------------------------------
\section{Two-Phase Design}
%---------------------------------------------------------------------------

The service is designed so that moving from Phase~1 to Phase~2 only replaces the data-source
layer. The DB schema, controllers, security stack, and the API the frontend calls are identical
in both phases.

\subsection{Phase 1 --- Foundation (current scope)}

\textbf{Goal:} get the Alert page reading from the microservice's own database instead
of the backend directly.

\begin{table}[h]
\small
\begin{tabularx}{\textwidth}{l X}
\toprule
\textbf{Component} & \textbf{Description} \\
\midrule
Node.js + Express service & Standalone HTTP server at its own URL \\
Alert MySQL database & Own database: \texttt{alerts}, \texttt{alert\_counts}, \texttt{service\_keys} \\
Sync middleware & Fetches from \texttt{UPSTREAM\_BASE\_URL}/Laravel, upserts into alert DB \\
Security stack & Helmet, CORS, rate limiter, JWT validation, input validation \\
REST API & 5 endpoints: counts + 4 paginated lists \\
Frontend service file & \texttt{alertService.js} --- the only frontend file that knows the microservice URL \\
\bottomrule
\end{tabularx}
\end{table}

\textbf{What the frontend Alert page gains:}
\begin{itemize}
  \item One clean URL for all alert data
  \item Consistent camelCase response shapes
  \item Rejection reasons included directly in the rejected list --- no secondary API calls
  \item A security-hardened access layer for alert traffic
  \item A foundation ready for Phase~2 with no frontend changes required
\end{itemize}

\subsection{Phase 2 --- Event-driven (future scope)}

\textbf{Goal:} remove the polling dependency on the upstream backend; producers publish
alert events directly to the service.

\begin{table}[h]
\small
\begin{tabularx}{\textwidth}{l X X}
\toprule
\textbf{What} & \textbf{Phase 1} & \textbf{Phase 2} \\
\midrule
Data source & Sync middleware polls upstream & Producers POST events here \\
Laravel role & Source of data & Publisher of events \\
New services & Cannot contribute alerts & Publish with a service API key \\
DB schema & Designed for Phase~2 & Unchanged \\
Frontend API & --- & Unchanged \\
Security & User JWT + CORS + rate limit & Same + API key auth for producers \\
\bottomrule
\end{tabularx}
\end{table}

The DB schema is designed for Phase~2 from day one so no migration is needed at the switch.

%---------------------------------------------------------------------------
\section{Security Architecture}
\label{sec:security}
%---------------------------------------------------------------------------

\begin{keybox}[Defense in depth]
Every request passes through a layered security stack. A request that fails any layer is
rejected immediately and never reaches the next layer.
\end{keybox}

\subsection{Middleware execution order}

\begin{enumerate}
  \item \textbf{Helmet} --- set security headers on every response (before anything else)
  \item \textbf{CORS} --- reject requests from unknown origins
  \item \textbf{Rate Limiter} --- reject if request rate exceeds threshold
  \item \textbf{Auth / Service Auth} --- validate JWT (user routes) or API key (producer routes)
  \item \textbf{Input Validation} --- reject malformed or out-of-range parameters
  \item \textbf{Controller} --- business logic, calls service layer
  \item \textbf{Error Handler} --- catch all thrown errors, return safe response
\end{enumerate}

\subsection{Transport security}

All traffic in production goes over \textbf{HTTPS only}. TLS termination is handled at the
infrastructure layer (nginx-proxy or load balancer), consistent with the existing EMIS stack.
The Express app runs HTTP behind the TLS terminator.

\subsection{Security headers (Helmet)}

\texttt{helmet} is applied as the \textbf{first middleware} --- before CORS, auth, or anything
else. It sets protective HTTP response headers on every response.

\begin{table}[h]
\small
\begin{tabularx}{\textwidth}{l l X}
\toprule
\textbf{Header} & \textbf{Value} & \textbf{Purpose} \\
\midrule
\texttt{X-Content-Type-Options} & \texttt{nosniff} & Prevent MIME-type sniffing \\
\texttt{X-Frame-Options} & \texttt{DENY} & Prevent clickjacking \\
\texttt{Strict-Transport-Security} & \texttt{max-age=31536000} & Enforce HTTPS for 1 year \\
\texttt{Content-Security-Policy} & \texttt{default-src 'none'} & Block all content embedding \\
\texttt{X-Powered-By} & (removed) & Do not expose the framework name \\
\texttt{Referrer-Policy} & \texttt{no-referrer} & Do not leak URL in Referer header \\
\bottomrule
\end{tabularx}
\end{table}

\subsection{CORS}

Strict, explicit policy --- no wildcard.

\begin{table}[h]
\small
\begin{tabularx}{\textwidth}{l X}
\toprule
\textbf{Setting} & \textbf{Value} \\
\midrule
\texttt{origin} & \texttt{CORS\_ORIGIN} env var only (exact frontend URL) \\
\texttt{methods} & \texttt{GET, OPTIONS} (user routes only --- no POST/DELETE from browser) \\
\texttt{allowedHeaders} & \texttt{Authorization, Content-Type} only \\
\texttt{credentials} & \texttt{true} \\
\bottomrule
\end{tabularx}
\end{table}

\subsection{Rate limiting}

Three tiers of rate limiting using \texttt{express-rate-limit}, applied per IP address.
Rate limit headers are returned in every response so clients can self-throttle.

\begin{table}[h]
\small
\begin{tabularx}{\textwidth}{l l l l X}
\toprule
\textbf{Tier} & \textbf{Applied to} & \textbf{Window} & \textbf{Max requests} & \textbf{On breach} \\
\midrule
Global & All routes & 15 min & 500 & \texttt{429 TOO\_MANY\_REQUESTS} \\
Alert reads & \texttt{GET /api/alerts/*} & 1 min & 60 & \texttt{429 TOO\_MANY\_REQUESTS} \\
Health check & \texttt{GET /health} & 1 min & 120 & \texttt{429 TOO\_MANY\_REQUESTS} \\
\bottomrule
\end{tabularx}
\end{table}

\subsection{JWT validation (user auth)}

\begin{warnbox}[Key change from original design]
The alert service \textbf{validates JWTs itself} using the WSO2 IS JWKS endpoint. It does not
rely solely on the upstream backend to be the gatekeeper. A request with a fake, expired, or
malformed token is rejected before any DB or upstream sync call is made.
\end{warnbox}

Token validation now occurs at \textbf{three points}:
\begin{enumerate}
  \item \textbf{Alert Service} (new) --- signature + claims verified before anything else happens
  \item \textbf{APIM} (when deployed) --- full OAuth token introspection
  \item \textbf{Laravel backend} --- JwtGuard JWKS validation
\end{enumerate}

The \texttt{auth} middleware:
\begin{enumerate}
  \item Extracts the Bearer token. Missing or malformed $\Rightarrow$ \texttt{401} immediately.
  \item Fetches the JWKS from \texttt{WSO2\_JWKS\_URI} (cached in memory for \texttt{JWKS\_CACHE\_TTL}
        seconds, default 3600 --- matching Laravel's behaviour).
  \item Verifies the token signature using the matching public key.
  \item Verifies \texttt{exp} (not expired), \texttt{iss} (matches \texttt{WSO2\_ISSUER}),
        and \texttt{aud} (if \texttt{WSO2\_AUDIENCE} is set).
  \item If any check fails $\Rightarrow$ \texttt{401} immediately. No sync call is made.
        No DB is touched.
  \item If all checks pass $\Rightarrow$ decoded payload attached to \texttt{req.jwtPayload}.
\end{enumerate}

\textbf{Libraries:} \texttt{jsonwebtoken} for verification; \texttt{jwks-rsa} for JWKS
fetching and automatic key rotation handling.

\subsection{Service-to-service auth (Phase 2)}

When producers publish alert events they use a dedicated \textbf{API key}, not a user JWT.

\begin{itemize}
  \item Header: \texttt{X-Service-Key: \textless api-key\textgreater}
  \item Keys are generated as 64-character cryptographically random hex strings
  \item Only the \textbf{SHA-256 hash} of the key is stored in the \texttt{service\_keys} DB table ---
        the plain-text key is never stored anywhere
  \item A compromised key is revoked by setting \texttt{active = 0} in the table --- no redeploy needed
  \item This middleware is applied \textbf{only} to producer routes (\texttt{POST /api/events/*});
        user-facing read routes never see it
\end{itemize}

\subsection{Input validation}

All query parameters are validated and sanitized using \texttt{express-validator} before the
controller is reached. A request with invalid parameters is rejected with \texttt{400 BAD\_REQUEST}.

\begin{table}[h]
\small
\begin{tabularx}{\textwidth}{l l X}
\toprule
\textbf{Parameter} & \textbf{Rule} & \textbf{On failure} \\
\midrule
\texttt{page} & Integer, min 1, max 10000 & \texttt{400 BAD\_REQUEST} \\
\texttt{per\_page} & Integer, min 1, max 100 & \texttt{400 BAD\_REQUEST} \\
\bottomrule
\end{tabularx}
\end{table}

Extra query parameters not in the allowed list are stripped silently before the request
reaches the controller.

\subsection{Database security}

\begin{itemize}
  \item \textbf{Restricted DB user:} \texttt{DB\_USER} is granted only \texttt{SELECT},
        \texttt{INSERT}, \texttt{UPDATE} on the alert tables. No \texttt{DROP}, no \texttt{GRANT}.
        A separate \texttt{DB\_MIGRATE\_USER} with broader privileges is used only at startup
        for migrations.
  \item \textbf{Parameterized queries only:} all DB queries use \texttt{mysql2} parameterized
        syntax. String concatenation into SQL is never used. This fully prevents SQL injection.
  \item \textbf{Sensitive data:} the \texttt{payload} JSON column stores only display data
        (names, school, appointment dates). No NIC numbers, addresses, or credentials.
\end{itemize}

\subsection{Error handling --- no information leakage}

\begin{itemize}
  \item In \texttt{production}, all errors return a generic message to the client. No stack
        traces, no DB query text, no internal details.
  \item Full error detail (message + stack + request ID) is written to server logs only.
  \item Upstream errors (backend, future APIM) are translated --- the client sees
        \texttt{UPSTREAM\_ERROR}, not the raw upstream response body.
  \item Error responses never include the values of request headers or query parameters.
\end{itemize}

\textbf{Consistent error response shape:}
\begin{quote}
\small\ttfamily
\{ "success": false, "error": \{ "code": "UNAUTHORIZED", "message": "Token validation failed." \} \}
\end{quote}

\subsection{Request logging}

The \texttt{requestLogger} middleware logs request ID, method, path, status code, response
time, and client IP. It \textbf{never} logs the value of the \texttt{Authorization} header,
the \texttt{X-Service-Key} header, any part of the JWT payload, or any request/response body.

\subsection{Dependency security}

\begin{itemize}
  \item \texttt{package-lock.json} is committed --- exact versions are pinned
  \item \texttt{npm audit} runs in CI/CD on every push; builds fail on high-severity findings
  \item No new dependency is added without reviewing its CVE history
\end{itemize}

\subsection{Container security}

\begin{itemize}
  \item Runs as a \textbf{non-root user} (UID 1000) inside the container
  \item \textbf{Multi-stage Docker build} --- only production dependencies in the final image,
        no dev tools, no build artifacts
  \item Uses the official \texttt{node:20-alpine} base image --- minimal attack surface
  \item \texttt{.dockerignore} excludes \texttt{.env}, \texttt{node\_modules}, test files,
        and documentation
\end{itemize}

%---------------------------------------------------------------------------
\section{Project Structure}
%---------------------------------------------------------------------------

The service lives at the repo root as a peer of \texttt{backend/} and \texttt{frontend/}:

\begin{quote}
\small\ttfamily
emis/ \\
\hspace*{1em}backend/\hspace{2.5cm} Laravel + FrankenPHP (existing) \\
\hspace*{1em}frontend/\hspace{2.4cm} React + Vite (existing) \\
\hspace*{1em}alert-service/\hspace{1.6cm} Node.js + Express + MySQL (\textbf{NEW}) \\
\hspace*{1em}infrastructure/ \\
\hspace*{1em}docs/
\end{quote}

Inside \texttt{alert-service/src/}:

\begin{quote}
\small\ttfamily
src/ \\
\hspace*{1em}app.js\hspace{3.7cm} Express app setup, middleware registration in order \\
\hspace*{1em}server.js\hspace{3.3cm} HTTP server bootstrap \\
\hspace*{1em}config/index.js\hspace{2cm} Validate env vars at startup; crash on missing required vars \\
\hspace*{1em}routes/ \\
\hspace*{2em}alertRoutes.js\hspace{1.8cm} GET /api/alerts/* (user-facing) \\
\hspace*{2em}eventRoutes.js\hspace{1.8cm} POST /api/events/* (Phase 2 producer routes, stubbed) \\
\hspace*{1em}controllers/ \\
\hspace*{2em}alertController.js\hspace{1.1cm} Extract validated params $\rightarrow$ call service $\rightarrow$ respond \\
\hspace*{1em}services/ \\
\hspace*{2em}alertService.js\hspace{1.5cm} Query alert DB, return shaped data \\
\hspace*{2em}syncService.js\hspace{1.5cm} Fetch from upstream, upsert into alert DB \\
\hspace*{1em}db/ \\
\hspace*{2em}connection.js\hspace{1.9cm} mysql2 connection pool \\
\hspace*{2em}migrations/ \\
\hspace*{3em}001\_create\_alerts.sql \\
\hspace*{3em}002\_create\_alert\_counts.sql \\
\hspace*{3em}003\_create\_service\_keys.sql \\
\hspace*{1em}middleware/ \\
\hspace*{2em}helmet.js\hspace{2.5cm} Security headers (first middleware registered) \\
\hspace*{2em}cors.js\hspace{2.7cm} CORS --- locked to CORS\_ORIGIN \\
\hspace*{2em}rateLimiter.js\hspace{1.7cm} Three-tier rate limiting \\
\hspace*{2em}auth.js\hspace{2.7cm} Full JWT validation (signature + claims) \\
\hspace*{2em}serviceAuth.js\hspace{1.8cm} Phase 2 API key validation for producers \\
\hspace*{2em}validate.js\hspace{2.2cm} Input sanitization (express-validator) \\
\hspace*{2em}errorHandler.js\hspace{1.5cm} Centralised safe error responses \\
\hspace*{2em}requestLogger.js\hspace{1.2cm} Structured logging, never logs tokens or keys
\end{quote}

\begin{table}[h]
\small
\begin{tabularx}{\textwidth}{l l X}
\toprule
\textbf{Layer} & \textbf{File} & \textbf{Responsibility} \\
\midrule
Security & \texttt{helmet.js} & Set HTTP security headers on every response \\
Security & \texttt{cors.js} & Reject requests from unknown origins \\
Security & \texttt{rateLimiter.js} & Three-tier rate limiting per IP \\
Auth & \texttt{auth.js} & Validate JWT signature + claims (JWKS-based) \\
Auth & \texttt{serviceAuth.js} & Validate producer API key (Phase 2) \\
Validation & \texttt{validate.js} & Sanitize and validate all query parameters \\
Route & \texttt{alertRoutes.js} & Map HTTP method + path to controller \\
Controller & \texttt{alertController.js} & Extract params, trigger sync, call service, respond \\
Service & \texttt{alertService.js} & Query the alert DB, return shaped data \\
Sync & \texttt{syncService.js} & Fetch from upstream, normalize, upsert into DB \\
DB & \texttt{db/connection.js} & Manage mysql2 connection pool \\
Error & \texttt{errorHandler.js} & Catch all errors; safe response; full detail in logs \\
Logging & \texttt{requestLogger.js} & Structured request logging (no secrets) \\
Config & \texttt{config/index.js} & Validate env vars at startup; fail fast on missing \\
\bottomrule
\end{tabularx}
\end{table}

%---------------------------------------------------------------------------
\section{Technology Stack}
%---------------------------------------------------------------------------

\begin{table}[h]
\small
\begin{tabularx}{\textwidth}{l l X}
\toprule
\textbf{Concern} & \textbf{Technology} & \textbf{Reason} \\
\midrule
Runtime & Node.js 20 LTS & Lightweight, fast non-blocking I/O \\
Framework & Express 4 & Minimal, well-understood, no magic \\
Document store & MongoDB Community & Self-hosted; flexible document model for variable alert payloads; no overseas dependency \\
DB ODM & \texttt{mongoose} & Schema definition, validation, and query abstraction for MongoDB \\
Event queue (Phase 2) & Redis & Self-hosted; Redis Streams for durable async event delivery from producers \\
HTTP client & \texttt{axios} & Same library used in the frontend \\
JWT validation & \texttt{jsonwebtoken} + \texttt{jwks-rsa} & Industry standard; same approach as Laravel's \texttt{firebase/php-jwt} \\
Security headers & \texttt{helmet} & One dependency, covers all OWASP header recommendations \\
Rate limiting & \texttt{express-rate-limit} & Lightweight, no external dependency \\
Input validation & \texttt{express-validator} & Declarative, well-maintained \\
Environment & \texttt{dotenv} & Standard; \texttt{.env.example} pattern already used across this repo \\
Logging & \texttt{morgan} + structured \texttt{console} & No external log service dependency in Phase 1 \\
Containerisation & Docker (multi-stage, non-root) & Consistent with \texttt{backend/} and \texttt{frontend/} \\
\bottomrule
\end{tabularx}
\end{table}

%---------------------------------------------------------------------------
\section{R\&D: Database Technology Selection}
%---------------------------------------------------------------------------

Before committing to a database engine, an evaluation was conducted to determine whether
a relational (SQL) or non-relational (NoSQL) database better fits the alert service.
This section documents the constraints, findings, and recommendation from that evaluation.

\subsection{Evaluation constraints}

Two hard constraints were identified for any database used in this service:

\begin{itemize}
  \item \textbf{Free and open source.} The service must not introduce a paid database licence.
  \item \textbf{Self-hosted only --- no overseas location.} As part of the EMIS system operated
        by the Ministry of Education, Sri Lanka, all data must remain on infrastructure inside
        the country. This is a data sovereignty requirement. It rules out all managed cloud
        database services (MongoDB Atlas, AWS DynamoDB, AWS DocumentDB, Azure Cosmos DB,
        Google Firestore, etc.) regardless of their technical merits --- they host data
        on overseas servers.
\end{itemize}

\subsection{Why NoSQL is a strong candidate}

Alert and notification data has characteristics that favour a document-oriented store:

\begin{itemize}
  \item \textbf{Variable payload per producer.} Each producer service may include different
        fields in its alert payload. A document model stores these naturally without requiring
        schema migrations when a new producer is added.
  \item \textbf{Messaging and task handling patterns.} The intended use of this service ---
        collecting events from multiple producers and delivering them to role-based inboxes ---
        is a pattern where NoSQL document and key-value stores are widely used in production.
  \item \textbf{No complex joins required.} Alert queries are simple: filter by type, status,
        and role. There are no multi-table joins. A relational schema gives no structural advantage.
  \item \textbf{Future event volume.} As more producer services come online in Phase~2,
        the volume of alert writes will grow. Document stores generally scale write throughput
        more cheaply than relational engines.
\end{itemize}

\subsection{Candidate evaluation}

Only free, self-hostable engines were considered.

\begin{table}[h]
\small
\begin{tabularx}{\textwidth}{l l l X l}
\toprule
\textbf{Engine} & \textbf{Type} & \textbf{Free} & \textbf{Strengths for this use case} & \textbf{Complexity} \\
\midrule
\textbf{MongoDB Community} & Document & Yes &
  Flexible document model; native JSON queries; well-known driver ecosystem;
  runs easily in Docker; straightforward for alert storage with variable payloads &
  Low \\[4pt]
\textbf{Redis} & Key-value + Streams & Yes &
  Industry standard for queuing, pub/sub, and task queues; Redis Streams maps
  directly to Phase~2 event publishing; very low latency &
  Low \\[4pt]
\textbf{Apache Cassandra} & Column-family & Yes &
  Excellent for high-throughput time-series events; linear horizontal scale &
  High (requires cluster) \\[4pt]
\textbf{Apache CouchDB} & Document & Yes &
  Document model; HTTP-native API &
  Low \\[4pt]
\textbf{MySQL 8} & Relational & Yes &
  Already running in EMIS; native JSON column covers flexible payload;
  ACID transactions; familiar to the ops team &
  Low \\
\bottomrule
\end{tabularx}
\end{table}

\subsection{Recommendation}

\begin{keybox}[Selected engines]
\textbf{MongoDB Community} for persistent alert document storage. \\
\textbf{Redis} for Phase~2 async event queuing and pub/sub (when producers publish directly).
\end{keybox}

\textbf{MongoDB Community} is recommended for the alert document store because:
\begin{itemize}
  \item The alert record is naturally a document --- envelope fields (\texttt{alertType},
        \texttt{status}, \texttt{peopleId}) are fixed; payload content varies per producer.
        MongoDB stores both without a JSON-column workaround.
  \item Self-hosted in Docker alongside the other EMIS services --- no overseas dependency,
        no licence cost. Replaces the \texttt{alerts} and \texttt{alert\_counts} MySQL tables
        with a single \texttt{alerts} collection.
  \item Simple to operate at EMIS's current scale. A single replica-set node is sufficient
        for Phase~1 and can be expanded without application changes when needed.
\end{itemize}

\textbf{Redis} is recommended as a \textbf{Phase~2 addition} for the event queue when producers
begin publishing alert events asynchronously. Redis Streams provides durable, ordered event
delivery with consumer groups --- suited to the producer/consumer model in Phase~2. Redis does
not replace MongoDB; it sits in front of it as the intake queue.

\subsection{Why not MySQL for this service}

MySQL is already proven in EMIS and is not a wrong choice. The deciding factors for moving to
MongoDB here are:
\begin{itemize}
  \item The \texttt{payload JSON} workaround in MySQL is functional but turns the DB into a
        hybrid that is neither fully relational nor fully document-oriented.
  \item Adding new producer alert types in Phase~2 requires no schema changes in MongoDB ---
        a new producer simply adds new fields to its payload document.
  \item Keeping the alert service on its own NoSQL engine reinforces its independence from the
        main EMIS MySQL instance and makes the service easier to relocate or replicate separately.
\end{itemize}

%---------------------------------------------------------------------------
\section{Database Design}
%---------------------------------------------------------------------------

The MongoDB database uses three collections. The schema is designed generically so any
future producer can contribute alerts without a migration.

\textbf{Driver:} \texttt{mongoose} (ODM for Node.js + MongoDB). Replaces \texttt{mysql2}
in the technology stack.

\subsection{\texttt{alerts} collection}

One document per alert record. All producer-specific display data lives inside the
\texttt{payload} sub-document --- its shape varies per producer and requires no schema change
when a new producer type is added.

\begin{table}[h]
\small
\begin{tabularx}{\textwidth}{l X}
\toprule
\textbf{Field} & \textbf{Description} \\
\midrule
\texttt{\_id} & MongoDB ObjectId (auto-generated) \\
\texttt{alertType} & Category: \texttt{pending\_verification}, \texttt{revised}, \texttt{pending\_confirmation}, \texttt{rejected} \\
\texttt{producer} & Which system created this alert. Phase~1: \texttt{emis-hrm}. Future: any service name \\
\texttt{peopleId} & The teacher or principal this alert relates to \\
\texttt{recipientRole} & Which role should see this alert; \texttt{null} = all authorized roles \\
\texttt{status} & \texttt{active} = needs attention, \texttt{resolved} = actioned, \texttt{archived} = hidden \\
\texttt{payload} & Sub-document --- all display data (name, school, appointment date, rejection reason). Shape varies per producer; no schema constraint. \\
\texttt{externalRef} & ID in the producer's system; used for upsert deduplication during Phase~1 sync \\
\texttt{syncedAt} & Last time this record was refreshed from upstream (Phase~1 only) \\
\texttt{createdAt} & Document creation timestamp (auto-managed by Mongoose timestamps) \\
\bottomrule
\end{tabularx}
\end{table}

Indexes: \texttt{alertType}, \texttt{peopleId}, \texttt{status}, \texttt{syncedAt}.
Unique compound index: \texttt{\{ alertType, externalRef \}} for safe upsert deduplication.

\subsection{\texttt{alertcounts} collection}

One document per alert category, updated on every sync cycle. Powers the overview tab cards
and tab badge numbers. The \texttt{alertType} field is the unique key --- at most four documents
in this collection.

\subsection{\texttt{servicekeys} collection}

Stores hashed API keys for Phase~2 producers. Only the SHA-256 hash is stored; the
plain-text key is never written to the database.

\begin{table}[h]
\small
\begin{tabularx}{\textwidth}{l X}
\toprule
\textbf{Field} & \textbf{Description} \\
\midrule
\texttt{\_id} & MongoDB ObjectId \\
\texttt{serviceName} & Human-readable producer service name (e.g.\ \texttt{emis-hrm}) \\
\texttt{keyHash} & SHA-256 hash (hex, 64 chars) of the API key --- never the plain-text key \\
\texttt{active} & \texttt{true} = valid, \texttt{false} = revoked. Revoke without deleting for audit trail \\
\texttt{createdAt} & Key creation timestamp \\
\texttt{lastUsedAt} & Last successful authentication timestamp \\
\bottomrule
\end{tabularx}
\end{table}

%---------------------------------------------------------------------------
\section{API Surface (Frontend Contract)}
%---------------------------------------------------------------------------

The frontend API contract is \textbf{identical between Phase 1 and Phase 2} --- only
the internal data source changes.

\begin{table}[h]
\small
\begin{tabularx}{\textwidth}{l X}
\toprule
\textbf{Endpoint} & \textbf{Purpose} \\
\midrule
\texttt{GET /health} & Public health check --- no auth; used by Docker health checks \\
\texttt{GET /api/alerts/counts} & Count per alert category; powers overview cards and tab badges \\
\texttt{GET /api/alerts/pending-verification} & Paginated list of profiles awaiting first-level verification \\
\texttt{GET /api/alerts/revised} & Paginated list of profiles revised and resubmitted after rejection \\
\texttt{GET /api/alerts/pending-confirmation} & Paginated list of profiles awaiting employer confirmation \\
\texttt{GET /api/alerts/rejected} & Paginated list of rejected profiles, including rejection reason \\
\bottomrule
\end{tabularx}
\end{table}

\textbf{Request convention:} all protected endpoints require \texttt{Authorization: Bearer \textless token\textgreater}.
Missing or invalid token $\Rightarrow$ \texttt{401} immediately.

\textbf{Response conventions:}
\begin{itemize}
  \item Success: \texttt{\{ "success": true, "data": \textless payload\textgreater \}}
  \item Error: \texttt{\{ "success": false, "error": \{ "code": "...", "message": "..." \} \}}
  \item All property names are \textbf{camelCase} --- upstream snake\_case is normalized before storage
  \item Pagination via \texttt{page} and \texttt{per\_page} query parameters (max \texttt{per\_page}: 100)
\end{itemize}

%---------------------------------------------------------------------------
\section{Data Flow}
%---------------------------------------------------------------------------

\subsection{Counts request (Primary tab)}

\begin{enumerate}
  \item React \texttt{Alert.jsx} mounts; \texttt{alertService.getAlertCounts()} is called
  \item \texttt{GET /api/alerts/counts} with \texttt{Authorization: Bearer \textless token\textgreater}
  \item Helmet applies security headers
  \item CORS check passes (frontend origin is allowed)
  \item Rate limiter: within limit, continue
  \item \texttt{auth} middleware: extract Bearer token $\rightarrow$ fetch JWKS (from cache) $\rightarrow$
        verify signature, \texttt{exp}, \texttt{iss} $\rightarrow$ valid; attach payload to request
  \item \texttt{alertController.getCounts()} called
  \item \texttt{syncService.syncCounts(token)}: calls \texttt{UPSTREAM\_BASE\_URL/alerts/counts}
        with Bearer token $\rightarrow$ upstream validates and returns counts $\rightarrow$ upserted
        into \texttt{alert\_counts} table
  \item \texttt{alertService.getCounts()}: \texttt{SELECT} from \texttt{alert\_counts}
  \item Controller sends \texttt{200} with normalized \texttt{\{ pendingVerification, revised, ... \}}
  \item React renders count cards
\end{enumerate}

\subsection{List request (any tab)}

\begin{enumerate}
  \item User clicks a tab; the list component calls, e.g.,
        \texttt{alertService.getPendingVerification(\{ page: 1, perPage: 20 \})}
  \item \texttt{GET /api/alerts/pending-verification?page=1\&per\_page=20} + Bearer token
  \item Helmet $\rightarrow$ CORS $\rightarrow$ rate limiter $\rightarrow$ JWT validation (all pass)
  \item \texttt{validate} middleware: \texttt{page=1} valid; \texttt{per\_page=20} valid
  \item \texttt{syncService.syncList(type, token, page, perPage)}: fetches paginated list from
        upstream $\rightarrow$ normalizes $\rightarrow$ upserts into \texttt{alerts} table
  \item \texttt{alertService.getList(type, page, perPage)}: parameterized \texttt{SELECT}
        from \texttt{alerts} with \texttt{LIMIT}/\texttt{OFFSET}
  \item Controller sends \texttt{200} with \texttt{\{ items: [...], pagination: \{...\} \}}
  \item Component sets state; passes data as props to \texttt{AlertTeacherList}
\end{enumerate}

%---------------------------------------------------------------------------
\section{Environment Configuration}
%---------------------------------------------------------------------------

\begin{table}[h]
\small
\begin{tabularx}{\textwidth}{l l X}
\toprule
\textbf{Variable} & \textbf{Required} & \textbf{Description} \\
\midrule
\texttt{PORT} & No (3001) & Port the Express server listens on \\
\texttt{NODE\_ENV} & No (development) & Controls error detail in responses \\
\texttt{UPSTREAM\_BASE\_URL} & Yes & Backend URL (Laravel direct now; APIM in $\sim$1 month) \\
\texttt{CORS\_ORIGIN} & Yes & Exact frontend origin; never a wildcard \\
\texttt{WSO2\_JWKS\_URI} & Yes & JWKS endpoint on WSO2 IS / Asgardeo \\
\texttt{WSO2\_ISSUER} & Yes & Expected \texttt{iss} claim in JWT \\
\texttt{WSO2\_AUDIENCE} & No & Expected \texttt{aud} claim; if set, tokens without it are rejected \\
\texttt{JWKS\_CACHE\_TTL} & No (3600) & Seconds to cache JWKS public keys in memory \\
\texttt{DB\_HOST} & Yes & MySQL host \\
\texttt{DB\_PORT} & No (3306) & MySQL port \\
\texttt{DB\_NAME} & Yes & Alert service database name (e.g.\ \texttt{emis\_alerts}) \\
\texttt{DB\_USER} & Yes & Restricted MySQL user: SELECT, INSERT, UPDATE only \\
\texttt{DB\_PASSWORD} & Yes & MySQL password \\
\texttt{DB\_MIGRATE\_USER} & Yes & Migration user with broader privileges; used at startup only \\
\texttt{DB\_MIGRATE\_PASSWORD} & Yes & Migration user password \\
\bottomrule
\end{tabularx}
\end{table}

%---------------------------------------------------------------------------
\section{Impact on the Frontend Alert Page}
%---------------------------------------------------------------------------

The Alert page (Inbox) continues to look and work exactly the same for users. Under the hood:

\begin{itemize}
  \item Each component now owns its data fetching and passes data down as props ---
        clean separation of data and display
  \item All alert API calls go through a new \texttt{alertService.js} file --- changing
        the data source in the future requires editing one file only
  \item \texttt{AlertTeacherList} becomes a pure display component --- no fetch logic,
        easier to test and reuse
  \item Rejection reason data is included directly in the list response from the microservice
        --- the existing secondary API call to fetch rejection comments is eliminated
\end{itemize}

\textbf{No visible change to users. Cleaner, safer code for developers.}

%---------------------------------------------------------------------------
\section{What This Enables in the Future}
%---------------------------------------------------------------------------

\begin{table}[h]
\small
\begin{tabularx}{\textwidth}{X X}
\toprule
\textbf{Future capability} & \textbf{What it requires} \\
\midrule
Any new EMIS service sends alerts & Add \texttt{POST /api/events/alert} with service API key auth \\
Real-time count updates & Add WebSocket / Server-Sent Events layer on top of existing DB \\
Email or SMS notifications & Alert service triggers on new alert records \\
Notification history / audit trail & Already in the DB from day one --- expose a history endpoint \\
Mobile app inbox & Same REST API, same contract, same auth \\
APIM integration & Change one env var (\texttt{UPSTREAM\_BASE\_URL}), redeploy \\
\bottomrule
\end{tabularx}
\end{table}

None of these require touching the main Laravel backend or the frontend Alert page structure.

%---------------------------------------------------------------------------
\section{Deployment}
%---------------------------------------------------------------------------

\subsection{Local development}

\begin{quote}
\small\ttfamily
cd alert-service \\
cp .env.example .env \quad\textit{\# fill in DB\_* and UPSTREAM\_BASE\_URL} \\
npm install \\
npm run dev \quad\quad\quad\quad\textit{\# nodemon, restarts on file changes}
\end{quote}

Service available at \texttt{http://localhost:3001}. Requires a local MySQL instance with
the \texttt{emis\_alerts} database created. Migrations run automatically on startup using
\texttt{DB\_MIGRATE\_USER}.

\subsection{Docker (local stack)}

Added to \texttt{infrastructure/docker-swarm/docker-compose.local.yml}:
\begin{itemize}
  \item \texttt{alert-service} container (Node.js app, non-root user)
  \item \texttt{alert-mysql} container (dedicated MySQL 8 instance)
\end{itemize}
Both join the existing Docker overlay network.

\subsection{Production}

Added to \texttt{infrastructure/docker-swarm/docker-stack.yml}. The production MySQL for
the alert service is a \textbf{separate instance} from the main EMIS MySQL --- enforcing
DB-level isolation between services. The \texttt{alert-service} container runs as UID 1000
(non-root).

\subsection{Pre-production security checklist}

\begin{itemize}
  \item[$\square$] \texttt{NODE\_ENV=production} is set
  \item[$\square$] \texttt{DB\_USER} is a restricted MySQL user (no \texttt{DROP}, no \texttt{GRANT})
  \item[$\square$] \texttt{CORS\_ORIGIN} is the exact production frontend URL --- no wildcard
  \item[$\square$] \texttt{WSO2\_JWKS\_URI} and \texttt{WSO2\_ISSUER} point to the production IS instance
  \item[$\square$] TLS termination is in place at the nginx-proxy or load balancer layer
  \item[$\square$] \texttt{npm audit} returns no high or critical vulnerabilities
  \item[$\square$] Docker image built with multi-stage Dockerfile (not dev mode)
  \item[$\square$] \texttt{.env} file is not committed and not baked into the Docker image
  \item[$\square$] \texttt{DB\_MIGRATE\_USER} credentials are rotated after initial migration
\end{itemize}

%---------------------------------------------------------------------------
\section{Summary}
%---------------------------------------------------------------------------

\begin{table}[h]
\small
\begin{tabularx}{\textwidth}{l X X}
\toprule
 & \textbf{Before} & \textbf{After} \\
\midrule
Alert data lives in & Laravel's shared database & Dedicated MongoDB alert service \\
Database engine & MySQL (shared) & MongoDB Community (self-hosted, no overseas) \\
Phase 2 event queue & N/A & Redis Streams (self-hosted) \\
Frontend talks to & Laravel backend directly & Alert microservice \\
JWT validation & Upstream only & Alert service + upstream (defense in depth) \\
Rate limiting & None for alert traffic & Three-tier per-IP limiting \\
Future services publish alerts & No --- must write to Laravel DB & Yes --- API key to alert service \\
Alert API consistency & 3 patterns across 5 files & One file, one contract \\
Notification foundation & Not present & In place from day one \\
APIM readiness & N/A & One env var change \\
\bottomrule
\end{tabularx}
\end{table}

\vspace{0.8cm}
\noindent\textit{End of proposal. \quad Branch: \texttt{be/feat/alert-service}}

\end{document}
